%=====================================================================
\chapter{Correctness: Loop Invariants}\label{ch:invariants}
%=====================================================================
\locator{1}{Part I --- Analysing an Algorithm}

\begin{outcomes}
By the end of this chapter you will be able to:
\begin{tight}
  \item \textbf{Explain} why testing cannot establish that an algorithm is
        correct, and \textbf{quantify} the input space it would have to cover.
  \item \textbf{State} a loop invariant for a given loop, and \textbf{name} the
        three obligations a proof must discharge.
  \item \textbf{Prove} insertion sort and binary search correct from their
        invariants, including the permutation clause.
  \item \textbf{Distinguish} partial correctness from total correctness, and
        \textbf{prove} termination by exhibiting a decreasing measure.
  \item \textbf{Diagnose} a loop whose invariant holds but which fails to
        terminate.
\end{tight}
\end{outcomes}

\begin{prereq}
\chref{role} for the vocabulary. Section~4.2 leans on mathematical induction;
Appendix~B has it if you want a refresher first.
\end{prereq}

\begin{keyterms}
correctness $\bullet$ partial correctness $\bullet$ total correctness $\bullet$
loop invariant $\bullet$ initialisation $\bullet$ maintenance $\bullet$
termination $\bullet$ precondition $\bullet$ postcondition $\bullet$
permutation $\bullet$ loop variant $\bullet$ decreasing measure $\bullet$
well-founded
\end{keyterms}

\section{Why Testing Is Not Enough}

A sorting routine takes any array of any length holding any values. That is an
infinite set of inputs, and a test suite is a finite list. The gap between them
is not a matter of diligence; it is a matter of cardinality.

\begin{worked}[Counting What Exhaustive Testing Would Cost]
Restrict the problem absurdly: arrays of exactly $n$ elements, each value drawn
from $0$ to $99$. How many distinct sorted inputs must a binary search be
tested against?

\smallskip
\textbf{Step 1.} Only the multiset of values matters, so the count is the number
of multisets of size $n$ from an alphabet of $m = 100$:
\[ \binom{m + n - 1}{n}. \]

\smallskip
\textbf{Step 2.} Evaluated by \texttt{code/ch04\_invariant.cpp}, and each array
must then be tried against all 100 possible keys:

\rowcolors{2}{white}{lightbg}
\begin{center}
\begin{tabular}{@{}rrr@{}}
\toprule
\rowcolor{primary!15}
$n$ & sorted arrays & $\times$ keys \\
\midrule
 4 & $4.42 \times 10^{6}$  & $4.4 \times 10^{8}$ \\
 8 & $3.26 \times 10^{11}$ & $3.3 \times 10^{13}$ \\
16 & $1.50 \times 10^{19}$ & $1.5 \times 10^{21}$ \\
20 & $2.46 \times 10^{22}$ & $2.5 \times 10^{24}$ \\
\bottomrule
\end{tabular}
\end{center}

\smallskip
\textbf{Step 3.} At a billion tests a second, the $n = 20$ row takes
$2.5 \times 10^{15}$ seconds --- about \textbf{80 million years}. And that is
for arrays of twenty elements holding two-digit numbers.

\smallskip
\textbf{Step 4 --- the conclusion.} Testing samples the input space; it cannot
cover it. A proof covers it, in half a page, for every $n$ and every alphabet
at once. The two activities answer different questions, and this chapter is
about the one testing cannot answer.
\end{worked}

\begin{conceptbox}[What testing is still for]
None of this says do not test. A proof establishes that the \emph{algorithm} is
correct; it says nothing about whether your \emph{program} implements that
algorithm, which is where most real defects live --- an off-by-one in a loop
bound, a wrong variable name, an integer overflow the model ignored.

\smallskip
Proof and testing catch disjoint classes of bug. The lab in this chapter
contains a defect that \emph{testing found in a third of its trials and a
careless proof would have missed}, and \chref{dptwo}'s lab contains the reverse.
Use both.
\end{conceptbox}

\section{The Loop Invariant}

\begin{definitionbox}[Loop Invariant]
A \term{loop invariant} is an assertion about the program's state that is true
immediately before the loop's first iteration and immediately before every
subsequent one.

\smallskip
A proof by loop invariant has \textbf{three obligations}:

\begin{tightnum}
  \item \textbf{Initialisation.} The invariant is true before the first
        iteration.
  \item \textbf{Maintenance.} If it is true before an iteration, it is true
        before the next one.
  \item \textbf{Termination.} The loop ends, and when it does the invariant ---
        together with the reason the loop stopped --- gives the result you
        wanted.
\end{tightnum}
\end{definitionbox}

\begin{conceptbox}[It is induction, wearing different clothes]
\rowcolors{2}{white}{lightbg}
\begin{longtable}{@{}L{4.0cm}L{4.0cm}L{6.0cm}@{}}
\toprule
\rowcolor{primary!15}
\textbf{Invariant proof} & \textbf{Induction} & \textbf{On what} \\
\midrule
\endhead
initialisation & base case & the loop has run 0 times \\
maintenance & inductive step & assume after $k$, show after $k+1$ \\
termination & --- & no counterpart: induction proves a property for all $k$,
  and says nothing about the loop stopping \\
\bottomrule
\end{longtable}

\smallskip
\textbf{The missing row is the whole point of Section~4.5.} Induction on the
iteration count gives you ``the invariant holds after every iteration that
happens''. Whether any particular iteration happens is a different question,
and nothing in the first two obligations answers it.
\end{conceptbox}

\begin{definitionbox}[Partial and Total Correctness]
An algorithm is \term{partially correct} if, \emph{whenever it halts}, its
output satisfies the postcondition.

\smallskip
It is \term{totally correct} if it is partially correct \emph{and} halts on
every input satisfying the precondition.

\smallskip
Initialisation and maintenance establish partial correctness. Termination is
what upgrades it, and Section~4.6 measures a program that has the first and not
the second.
\end{definitionbox}

\section{Insertion Sort, Proved}

\dsfig{\aoaalg{\algname{Insertion-Sort}$(A, n)$}{
\For{$j \gets 2$ \textbf{to} $n$}
  \State $\mathit{key} \gets A[j]$
    \Comment{the element to place}
  \State $i \gets j - 1$
  \While{$i > 0$ \textbf{and} $A[i] > \mathit{key}$}
    \State $A[i+1] \gets A[i]$
      \Comment{shift right}
    \State $i \gets i - 1$
  \EndWhile
  \State $A[i+1] \gets \mathit{key}$
\EndFor
}}{\algname{Insertion-Sort}. The outer loop's invariant concerns the prefix
$A[1\mathinner{\ldotp\ldotp}j-1]$; the inner loop's concerns where
$\mathit{key}$ has got to. Lines are numbered so the proof can cite them.}{insort}

The algorithm is indexed from 1, as the pseudocode throughout this book is, and
the C++ in the labs is indexed from 0. Appendix~A explains the translation; the
mathematics is cleaner from 1 and the code is cleaner from 0, and pretending
otherwise helps nobody.

\begin{theorem}[Insertion sort is correct]\label{thm:insort}
Given an array $A[1 \mathinner{\ldotp\ldotp} n]$ of $n \ge 1$ numbers,
\algname{Insertion-Sort} terminates with $A$ holding a permutation of its
original contents, in non-decreasing order.
\end{theorem}

\begin{proof}
Take as the invariant for the \textbf{for} loop:

\smallskip
\noindent\textit{At the start of each iteration with index $j$, the subarray
$A[1 \mathinner{\ldotp\ldotp} j-1]$ consists of the elements originally in
$A[1 \mathinner{\ldotp\ldotp} j-1]$, in non-decreasing order.}

\smallskip
\textbf{Initialisation.} Before the first iteration $j = 2$, so the subarray is
$A[1 \mathinner{\ldotp\ldotp} 1]$: one element, trivially sorted, and trivially
the element originally there.

\smallskip
\textbf{Maintenance.} Assume the invariant before the iteration with index $j$.
The inner \textbf{while} loop shifts $A[i]$ to $A[i+1]$ for decreasing $i$ while
$A[i] > \mathit{key}$, maintaining its own invariant: \emph{the entries
$A[i+2 \mathinner{\ldotp\ldotp} j]$ are those originally in
$A[i+1 \mathinner{\ldotp\ldotp} j-1]$, all of them $> \mathit{key}$, and
position $A[i+1]$ is free.} The loop exits when $i = 0$ or
$A[i] \le \mathit{key}$; in either case every entry left of $i+1$ is
$\le \mathit{key}$ and every entry right of it is $> \mathit{key}$, so writing
$\mathit{key}$ into the free slot $A[i+1]$ leaves
$A[1 \mathinner{\ldotp\ldotp} j]$ sorted. No element was overwritten without
first being copied, and $\mathit{key}$ was saved before the shifting began, so
the subarray still holds exactly the original elements. Incrementing $j$ makes
this the invariant for the next iteration.

\smallskip
\textbf{Termination.} The loop variable $j$ increases by one each iteration and
the loop ends when $j = n+1$. Substituting into the invariant: $A[1
\mathinner{\ldotp\ldotp} n]$ consists of the elements originally in $A[1
\mathinner{\ldotp\ldotp} n]$, in non-decreasing order --- which is the
postcondition. Termination itself is immediate because $j$ strictly increases
towards the fixed bound $n+1$, and the inner loop terminates because $i$
strictly decreases towards the guard $i > 0$.
\end{proof}

\begin{alertbox}[Do not drop the permutation clause]
It is tempting to state the invariant as ``$A[1\mathinner{\ldotp\ldotp}j-1]$ is
sorted'' and save eleven words. Then consider this implementation:
\begin{center}
\emph{overwrite every element of $A$ with zero.}
\end{center}
Its output is sorted. It satisfies the shortened invariant at every iteration
and terminates. It is not a sorting algorithm.

\smallskip
``Sorted'' is only half the postcondition; ``a permutation of the input'' is the
other half, and a proof that omits it proves the wrong theorem. This is the
single most common error in student invariant proofs.
\end{alertbox}

\section{Binary Search, Where the Invariant Is About Exclusion}

Insertion sort's invariant describes what has been \emph{built}. The other
common shape describes what has been \emph{ruled out}.

\dsfig{\aoaalg{\algname{Binary-Search}$(A, n, x)$}{
\State $\mathit{lo} \gets 1$;\quad $\mathit{hi} \gets n$
\While{$\mathit{lo} \le \mathit{hi}$}
  \State $\mathit{mid} \gets \lfloor (\mathit{lo} + \mathit{hi})/2 \rfloor$
  \If{$A[\mathit{mid}] = x$}
    \State \Return $\mathit{mid}$
  \ElsIf{$A[\mathit{mid}] < x$}
    \State $\mathit{lo} \gets \mathit{mid} + 1$
      \Comment{left half cannot hold $x$}
  \Else
    \State $\mathit{hi} \gets \mathit{mid} - 1$
      \Comment{right half cannot hold $x$}
  \EndIf
\EndWhile
\State \Return \algname{not-found}
}}{\algname{Binary-Search} on a sorted array. The $+1$ on line~7 and the $-1$ on
line~9 are what make the search range strictly smaller every iteration ---
Section~4.5 is about what happens when one of them is missing.}{binsearch}

\begin{theorem}[Binary search is correct]\label{thm:binsearch}
Let $A[1 \mathinner{\ldotp\ldotp} n]$ be sorted in non-decreasing order.
\algname{Binary-Search} returns an index $\mathit{mid}$ with
$A[\mathit{mid}] = x$ if $x$ occurs in $A$, and \algname{not-found} otherwise.
\end{theorem}

\begin{proof}
Invariant: \emph{if $x$ occurs anywhere in $A$, then it occurs in
$A[\mathit{lo} \mathinner{\ldotp\ldotp} \mathit{hi}]$.}

\smallskip
\textbf{Initialisation.} $\mathit{lo} = 1$ and $\mathit{hi} = n$, so the range
is all of $A$ and the implication is trivial.

\smallskip
\textbf{Maintenance.} Suppose the invariant holds and the loop body runs. If
$A[\mathit{mid}] = x$ the procedure returns a correct index and there is nothing
to maintain. If $A[\mathit{mid}] < x$ then, because $A$ is sorted, every index
$k \le \mathit{mid}$ has $A[k] \le A[\mathit{mid}] < x$, so no such position
holds $x$; discarding them by setting $\mathit{lo} \gets \mathit{mid}+1$
therefore cannot discard an occurrence of $x$, and the invariant survives. The
case $A[\mathit{mid}] > x$ is symmetric.

\smallskip
\textbf{Termination.} Each iteration either returns, or replaces
$\mathit{hi}-\mathit{lo}$ by a strictly smaller value: line~7 raises
$\mathit{lo}$ past $\mathit{mid} \ge \mathit{lo}$, and line~9 lowers
$\mathit{hi}$ below $\mathit{mid} \le \mathit{hi}$. Since
$\mathit{hi}-\mathit{lo}$ is an integer that strictly decreases and the loop
runs only while $\mathit{hi} - \mathit{lo} \ge 0$, the loop performs finitely
many iterations. If it exits, then $\mathit{lo} > \mathit{hi}$, so the range
$A[\mathit{lo} \mathinner{\ldotp\ldotp} \mathit{hi}]$ is empty; by the invariant
$x$ occurs in $A$ only if it occurs in an empty range, so it does not occur at
all, and \algname{not-found} is correct.
\end{proof}

\begin{corollary}[Binary search is logarithmic]
Binary search performs at most $\lfloor \lg n \rfloor + 1$ iterations, since
$\mathit{hi}-\mathit{lo}+1$ at least halves each time and the loop stops when
it reaches $0$.
\end{corollary}

\section{Termination: The Obligation People Skip}

The termination argument above did real work, and it is worth isolating what
kind of work.

\begin{definitionbox}[Loop Variant]
A \term{loop variant}, or decreasing measure, is a quantity $\mu$ computed from
the program state such that

\begin{tight}
  \item $\mu$ takes values in a \term{well-founded} set --- for our purposes,
        the non-negative integers;
  \item every iteration strictly decreases $\mu$.
\end{tight}

\smallskip
Then the loop terminates, because a strictly decreasing sequence of
non-negative integers must be finite.
\end{definitionbox}

The variants used so far: $n+1-j$ for insertion sort's outer loop, $i$ for its
inner loop, and $\mathit{hi}-\mathit{lo}+1$ for binary search.

\dsfig{\aoaalg{\algname{Euclid-GCD}$(a, b)$}{
\Require $a \ge 0$, $b \ge 0$, not both zero
\While{$b \ne 0$}
  \State $r \gets a \bmod b$
  \State $a \gets b$
  \State $b \gets r$
\EndWhile
\State \Return $a$
}}{\algname{Euclid-GCD}. Here the invariant is an equality --- $\gcd(a,b)$ never
changes --- and the variant is simply $b$, whose decrease is what makes the
loop finite.}{euclid}

\begin{theorem}[Euclid's algorithm is totally correct]\label{thm:euclid}
\algname{Euclid-GCD}$(a,b)$ terminates and returns $\gcd(a,b)$.
\end{theorem}

\begin{proof}
\textbf{Invariant:} $\gcd(a, b)$ equals the greatest common divisor of the
original arguments, and not both of $a,b$ are zero.

\smallskip
\textbf{Initialisation.} Immediate --- $a$ and $b$ are the original arguments.

\smallskip
\textbf{Maintenance.} The body replaces $(a,b)$ by $(b,\, a \bmod b)$. Write
$a = qb + r$ with $r = a \bmod b$. Any $d$ dividing both $b$ and $r$ divides
$qb + r = a$; any $d$ dividing both $a$ and $b$ divides $a - qb = r$. So the
pairs $(a,b)$ and $(b,r)$ have the same set of common divisors, hence the same
greatest one, and the invariant is preserved.

\smallskip
\textbf{Termination.} Take $\mu = b$. Since $b \neq 0$ in the body,
$r = a \bmod b$ satisfies $0 \le r < b$, so the new value of $b$ is strictly
less than the old one and never negative. A strictly decreasing sequence of
non-negative integers is finite, so the loop ends --- with $b = 0$. Then
$\gcd(a,0) = a$, and by the invariant $a$ is the gcd of the original arguments,
which is what line~6 returns.
\end{proof}

\begin{alertbox}[Where the measure must live]
``Strictly decreasing'' is not enough on its own; the measure must decrease
within a \emph{well-founded} set. Over the non-negative reals it is not:
\[ x \gets x/2 \qquad \text{from } x = 1 \]
strictly decreases forever and never terminates. Over the non-negative integers
no such sequence exists, which is the entire force of the argument.

\smallskip
Use integers, and if your measure is naturally a real number, find an integer
that bounds it.
\end{alertbox}

\section{A Loop Whose Invariant Holds and Which Never Finishes}

Here is the case the three obligations exist to catch. Change one character of
\algname{Binary-Search}: on line~7 write $\mathit{lo} \gets \mathit{mid}$
instead of $\mathit{lo} \gets \mathit{mid}+1$.

\begin{conceptbox}[The invariant still holds --- and that is the problem]
Re-read the maintenance argument of Theorem~\ref{thm:binsearch}. It establishes
that positions $k \le \mathit{mid}$ cannot hold $x$, and therefore that
discarding them is safe. Setting $\mathit{lo} \gets \mathit{mid}$ discards
\emph{fewer} positions than that --- it keeps $\mathit{mid}$ itself. Keeping a
position cannot lose an occurrence of $x$, so:

\smallskip
\textbf{Initialisation still holds. Maintenance still holds.} The modified
program is \emph{partially correct}: if it answers, the answer is right.

\smallskip
\textbf{Termination fails.} When $\mathit{hi} = \mathit{lo}+1$ we get
$\mathit{mid} = \lfloor (2\mathit{lo}+1)/2 \rfloor = \mathit{lo}$, so
$\mathit{lo} \gets \mathit{mid}$ leaves $\mathit{lo}$ and $\mathit{hi}$
unchanged and the next iteration is identical. The measure
$\mathit{hi}-\mathit{lo}$ no longer strictly decreases, and the loop spins
forever.
\end{conceptbox}

\begin{worked}[Measured: 132 Million Iterations, Zero Wrong Answers]
\texttt{code/ch04\_invariant.cpp} runs the broken search on a 64-element sorted
array, 10{,}000 trials with keys known to be present and 10{,}000 with keys
known to be absent, and \emph{audits the invariant on every single iteration}.
A run that exceeds 10{,}000 iterations is recorded as non-terminating.

\rowcolors{2}{white}{lightbg}
\begin{center}
\begin{tabular}{@{}lrr@{}}
\toprule
\rowcolor{primary!15}
10{,}000 trials each & non-termination & wrong answer \\
\midrule
keys always present & 3{,}228 & 0 \\
keys always absent  & 10{,}000 & 0 \\
\midrule
\textbf{iterations audited} & \multicolumn{2}{r}{\textbf{132{,}311{,}650}} \\
\textbf{invariant violations} & \multicolumn{2}{r}{\textbf{0}} \\
\bottomrule
\end{tabular}
\end{center}

\smallskip
\textbf{Read the second column first.} In 20{,}000 trials the broken search
returned a wrong answer \textbf{not once}. That is not luck; it is precisely
what the invariant guarantees, and the audit confirms the invariant never
broke across 132 million iterations.

\smallskip
\textbf{Now read the first column.} 13{,}228 of the 20{,}000 runs never
finished. The program is partially correct and totally useless.

\smallskip
\textbf{The distribution matters for how fast you find it.} Keys drawn from
inside the array hang about a third of the time; keys known to be absent hang
every time. A test suite built only from present keys would still have caught
this --- in roughly three trials --- which is worth saying plainly, because the
tidier story would be that testing misses it. Testing found this bug easily.

\smallskip
\textbf{What testing could not have told you is \emph{why}}, or that the
defect is confined entirely to termination with partial correctness intact. And
a proof that checked only initialisation and maintenance --- which is what most
students write --- would have certified this program as correct.
\end{worked}

\begin{pitfall}
\begin{tight}
  \item \textbf{Omitting the permutation clause} from a sorting invariant. The
        all-zeros ``sort'' satisfies the weakened version.
  \item \textbf{Stating an invariant that mentions no loop variable.} ``$A$ is
        an array of numbers'' is true throughout and proves nothing. A useful
        invariant relates the work done so far to the loop index.
  \item \textbf{Checking the invariant at the wrong moment.} It is asserted
        \emph{before} each iteration, not in the middle of the body, where it is
        routinely false --- insertion sort's prefix is genuinely unsorted while
        the shifting is in progress.
  \item \textbf{Treating termination as a formality.} Section~4.6 is a program
        that passes the other two obligations and runs forever.
  \item \textbf{A measure that decreases but not in a well-founded set.}
        $x \gets x/2$ over the reals.
  \item \textbf{Proving the loop correct and forgetting the exit condition.}
        The invariant alone is not the postcondition; you need the invariant
        \emph{and} the reason the loop stopped, which is why every termination
        paragraph above substitutes the final index back in.
  \item \textbf{Assuming a proved algorithm gives a correct program.} The
        proof is about the algorithm; overflow, off-by-one and typing errors
        live below it.
\end{tight}
\end{pitfall}

\begin{lab}[A Bug That Is Never Wrong]
Reproduces the table in Section~4.6. Expect twenty-five minutes; part C audits
132 million iterations and takes about a minute.

\begin{lstlisting}[language=C++]
// ch04_invariant.cpp -- why the THIRD obligation is not decoration.
// Build: cl /O2 /EHsc /std:c++17 ch04_invariant.cpp  (see Appendix A)

// --- 1. The correct version, and the invariant it maintains --------
// Invariant: if x occurs in a, then it occurs in a[lo..hi].
//   Initialisation: lo = 0, hi = n-1, so the range is all of a.
//   Maintenance:    a is sorted, so discarding the half that cannot
//                   contain x preserves the invariant.
//   Termination:    hi - lo strictly decreases and is bounded below.
static int bsearch_ok(const std::vector<int>& a, int x) {
    int lo = 0, hi = (int)a.size() - 1;
    while (lo <= hi) {
        int mid = lo + (hi - lo) / 2;
        if (a[mid] == x) return mid;
        if (a[mid] < x) lo = mid + 1;   // the correct step
        else            hi = mid - 1;
    }
    return -1;
}

// --- 2. The same search with one character changed -----------------
// "lo = mid" still discards nothing that could hold x, so the
// invariant SURVIVES. But when hi == lo + 1 and a[mid] < x, mid
// computes to lo, so lo = mid leaves lo and hi UNCHANGED and the loop
// makes no progress. It is not wrong; it is non-terminating.
static int bsearch_bug(const std::vector<int>& a, int x,
                       bool* inv_ever_broken, long long* iters) {
    int lo = 0, hi = (int)a.size() - 1;
    for (int step = 0; step < CAP; ++step) {
        if (lo > hi) return -1;
        int mid = lo + (hi - lo) / 2;

        // Audit the invariant itself, every single iteration: if x is
        // present in the whole array, is it still inside [lo, hi]?
        bool present = std::binary_search(a.begin(), a.end(), x);
        if (present) {
            bool in_window = false;
            for (int i = lo; i <= hi; ++i)
                if (a[i] == x) { in_window = true; break; }
            if (!in_window) *inv_ever_broken = true;
        }
        ++*iters;

        if (a[mid] == x) return mid;
        if (a[mid] < x) lo = mid;       // <-- the one-character bug
        else            hi = mid - 1;
    }
    return -2;                          // stands in for "forever"
}

// --- 3. Ten thousand tests with keys that are present --------------
// --- 4. The same ten thousand tests, keys absent -------------------
// The only change between parts 3 and 4 is the distribution the keys
// are drawn from, and it changes the hang rate from 32% to 100%.

// --- 5. And the invariant? It never broke once. --------------------
// inv_ever_broken is set by the audit in part 2 and is never set.
// Printing that zero beside the 13,228 hangs IS the experiment: the
// two numbers together say "partially correct, and useless".
printf("   iterations audited : %lld\n", iters);
printf("   invariant violated : %s\n", broken ? "YES" : "NO");
printf("   wrong answers, A+B : %d\n", wrong_present + wrong_absent);

// --- 6. Why proving beats testing, counted -------------------------
// Exhaustive testing means every sorted array and every key. The
// number of sorted arrays of length n over an alphabet of size m is
// the multiset coefficient C(m+n-1, n).
for (int n : {4, 8, 16, 20}) {
    const int m = 100;
    double c = 1.0;
    for (int i = 1; i <= n; ++i) c *= (double)(m + n - i) / i;
    printf("   n=%2d, values 0..99 : %.3g sorted arrays"
           " x %d keys\n", n, c, m);
}
\end{lstlisting}

\textbf{What to notice.} The \texttt{invariant violated} line prints
\texttt{NO}, and it is the most important line in the output: the audit is not
failing to find a violation because it is badly written, it is failing because
there is none to find.

\smallskip
\textbf{Then try to break it.} Change the bug from \texttt{lo = mid} to
\texttt{lo = mid + 2}. Now predict, before running it, which of the three
obligations fails --- and which column of the table changes. This mutation
terminates and gives wrong answers, the exact mirror image of the one in the
book, and it is the case an invariant proof catches immediately.
\end{lab}

\begin{checkpoint}
\begin{tightnum}
  \item State a loop invariant for a loop that computes $\sum_{i=1}^{n} A[i]$
        into a variable \textit{total}, and give the measure that proves it
        terminates.
  \item A classmate's invariant for selection sort reads: ``after each
        iteration, $A[1\mathinner{\ldotp\ldotp}k]$ is sorted''. Give an
        incorrect program that satisfies it, and repair the invariant.
  \item Explain, in terms of the three obligations, how a program can be
        partially correct but not totally correct, and give the measure whose
        absence causes it.
\end{tightnum}
\end{checkpoint}

\begin{chaptersummary}
\begin{tight}
  \item Testing samples an infinite input space. For sorted arrays of 20
        two-digit numbers, exhaustive testing needs $2.5 \times 10^{24}$ cases
        --- about 80 million years. A proof covers every $n$ at once.
  \item A loop invariant proof has three obligations: initialisation,
        maintenance, termination. The first two are induction; the third has no
        counterpart in induction and is where the real content often lies.
  \item Initialisation plus maintenance give \emph{partial} correctness --- if
        it halts, it is right. Termination upgrades that to total correctness.
  \item Insertion sort's invariant must say the prefix holds \emph{the elements
        originally there}, in order. Drop the permutation clause and the
        all-zeros ``sort'' satisfies it.
  \item Binary search's invariant is about exclusion: if $x$ is in $A$ at all,
        it is in $A[\mathit{lo}\mathinner{\ldotp\ldotp}\mathit{hi}]$.
  \item Termination is proved by a loop variant: an integer-valued measure that
        strictly decreases. Over the reals the argument fails --- $x \gets x/2$.
  \item Measured: a binary search with $\mathit{lo} \gets \mathit{mid}$ gave
        \textbf{zero} wrong answers in 20{,}000 trials and \textbf{zero}
        invariant violations across 132{,}311{,}650 audited iterations --- and
        failed to terminate in 13{,}228 of those trials. Two obligations out of
        three is not correctness.
\end{tight}
\end{chaptersummary}

\begin{reviewq}
  \item \textbf{(MCQ)} The three obligations of a loop invariant proof are:
        (a)~initialisation, iteration, conclusion; (b)~initialisation,
        maintenance, termination; (c)~precondition, invariant, postcondition;
        (d)~base case, inductive step, generalisation.
  \item \textbf{(MCQ)} Initialisation and maintenance alone establish:
        (a)~total correctness; (b)~partial correctness; (c)~termination;
        (d)~a running-time bound.
  \item \textbf{(Short)} Explain why a loop invariant proof is a proof by
        induction, and identify the one obligation that has no inductive
        counterpart.
  \item \textbf{(Short)} Why must a sorting algorithm's invariant mention that
        the sorted prefix contains the original elements? Give the program that
        makes the point.
  \item \textbf{(Short)} State the two conditions a loop variant must satisfy,
        and give an example of a measure satisfying the first but not the
        second.
  \item \textbf{(Applied)} Prove that the following loop computes $a^{n}$ for
        integer $n \ge 0$, discharging all three obligations:
        \textit{$r \gets 1$; $k \gets n$; while $k > 0$ do $r \gets r \cdot a$;
        $k \gets k - 1$}.
  \item \textbf{(Applied)} The measured experiment found a defect that is never
        wrong and sometimes never finishes. Construct the mirror-image defect
        --- one that always terminates and is sometimes wrong --- in binary
        search, say which obligation it violates, and describe the test that
        would find it.
\end{reviewq}
